\documentclass[12pt,titlepage]{article}
\usepackage[margin=1in]{geometry}
\usepackage{booktabs,tabularx,longtable,amsmath}
\setlength{\emergencystretch}{3em}
\input{../Common.text}
\newcommand{\case}[3]{\subsubsection{#1: #2}\label{test:#1}\noindent\textbf{Requirements:} #3.\par}
\begin{document}
\title{System Verification and Validation Plan for \progname}
\author{\authname}
\date{November 2, 2026}
\maketitle
\pagenumbering{roman}
\section*{Revision History}

\noindent\begin{tabularx}{\textwidth}{llX}
\toprule Date & Version & Notes\\\midrule
2026-11-02 & 0.1 & Initial cross-model KV prediction plan; proposed procedures, pending team review and setup decisions.\\\bottomrule
\end{tabularx}
\tableofcontents
\newpage
\pagenumbering{arabic}
\section{Symbols, Abbreviations, and Acronyms}

\noindent\begin{tabularx}{\textwidth}{lX}
\toprule Term & Meaning\\\midrule
KV & Attention key and value tensors.\\
SRS & Software Requirements Specification; requirement IDs below refer to the Volere SRS.\\
TTFT & Time from request arrival to first returned token.\\
p95 & 95th percentile.\\
Native & Selected target model using the serving engine without Tandem prediction.\\
Setup & Exact model pair, predictor, alignment, precision, software and hardware configuration.\\\bottomrule
\end{tabularx}

\section{General Information}
\subsection{Summary}
Tandem predicts target-model KV from source-model KV for shared conversation history. After validating the predicted cache, it loads that cache, prefills the remaining target tokens and decodes with the router-selected target model. This plan tests prediction, output consistency, serving integration and measured benefit on one server. It describes future work; no test result or setup approval is asserted here.

\subsection{Objectives}
The priorities are to preserve the selected model's behaviour and request isolation, demonstrate reduced target-prefix computation and net latency, and retain safe native processing. KV error and downstream output quality are evaluated separately. A small tensor error does not establish acceptable generated output.

Base-model training, router development, multi-server execution and universal correctness across all models/prompts are outside scope. Third-party components are tested at Tandem's integration boundaries. Their entire implementations are not reverified. Performance conclusions apply only to the recorded workloads and supported setups.

\subsection{Custom Documents}
The proposed Machine Learning Model Report should record model-pair selection, predictor architecture, alignment, fitting procedure, data splits and limitations. The existing proposal's statement that training does not apply needs review because predictor fitting is now in scope. Instructor approval remains pending. The proposed Performance Report will contain timing distributions, throughput, resource measurements, uncertainty and negative results; the V\&V Report will use this evidence when assessing requirements.

\subsection{Relevant Documentation}
The \href{../SRS-Volere/SRS.tex}{Volere SRS} supplies requirement IDs, recovery rules and draft acceptance criteria. The \href{../ProblemStatementAndGoals/ProblemStatement.tex}{Problem Statement} supplies project goals and scope. The \href{../DevelopmentPlan/DevelopmentPlan.tex}{Development Plan} supplies review roles and proposed Python tooling; its tree-attention and KV-Promise descriptions require alignment with cross-model prediction. The Hazard Analysis will be reviewed for cache isolation and lifecycle risks. MG/MIS documents will supply module boundaries for the later unit-test plan. Reviewers will use \href{../Checklists/VnV-Checklist.tex}{the V\&V checklist}.

The \href{../../src/experiments/LOCAL_CACHE_RESULTS.md}{local cache replay results} provide preliminary same-model cache controls and planted-code fixtures. They do not qualify cross-model prediction or establish Linux GPU serving performance. Existing routing tests can support integration checks, but router development remains outside this plan's scope.

\section{Plan}
\subsection{Verification and Validation Team}
The team comprises Justin Kwinecki, Zayed Umer, William Pol, Manav Patel and Musab Mondal. Named assignments will be recorded before execution; the following responsibilities are proposed rather than assigned.

\noindent\begin{tabularx}{\textwidth}{lX}
\toprule Role & Responsibility\\\midrule
Deliverable lead & Coordinate scope, traceability and unresolved decisions.\\
Prediction test owner & Prepare disjoint datasets, offline reference KV and output checks.\\
Integration test owner & Test validation, loading, fallback, isolation and cleanup.\\
Measurement owner & Freeze manifests, run matched trials and report uncertainty.\\
Independent reviewer & Review cases and evidence without reviewing their own work alone.\\\bottomrule
\end{tabularx}

\subsection{SRS Verification}
A reviewer will walk through a valid switch, a rejected cache, a failed prediction and a cancelled shared-cache request. For each, check that inputs, expected outcomes and requirement IDs are unambiguous. Review OI1--OI3 before final tests, especially whether native repeatability supports the draft output tolerances. Record contradictions and decisions as repository issues; thresholds must not be changed to rescue final-test failures.

\subsection{Design Verification}
Once the design is available, trace source KV through alignment, prediction, validation, loading, suffix prefill and cleanup. Check cache ownership, versioning, host-granted sharing permissions, asynchronous readiness and the boundary before first output where one native restart is allowed. Reviewers will record findings and link fixes before qualifying a setup.

\subsection{Verification and Validation Plan Verification}
An independent team member will check each case for a reproducible fixture, observable oracle, requirement link and feasible resource budget. Peer/TA feedback will be tracked through issues. Deliberately wrong prefix positions, bypassed permission checks and omitted prediction timing should cause the corresponding checks to fail; these are proposed checks of the test harness, not claims that mutation tests already exist.

\subsection{Implementation Verification}
Use peer code review, static checks and the system cases below. Review tensor layouts, token accounting, error paths and cache ownership explicitly. Each executed case will retain a commit ID, setup manifest, fixture ID, result and evidence location. Classify outcomes as pass, fail, blocked or inconclusive; an unavailable GPU or missing predictor is not a pass.

\subsection{Automated Testing and Verification Tools}
Use Python with pytest for fixtures, parameterized checks and fault injection; Ruff for PEP~8-oriented linting/formatting; and pytest-cov for line/branch coverage of Tandem-owned Python code. Coverage reports identify untested validation and recovery branches, without treating a coverage percentage as proof of correctness. These are planned checks; the existing documentation CI does not establish that code CI is configured.

CPU checks should run on pull requests. GPU correctness and timing runs require the pinned supported environment and will run before setup qualification and release. Use a monotonic clock for request timings, synchronized GPU events for device stages where applicable, and framework/driver memory counters for direct memory measurements. Record synchronization and instrumentation overhead. CPU or local MPS results cannot substitute for supported GPU performance trials.

\subsection{Software Validation}
Ask an independent technical user to install the supported setup, run a model-switch example and identify whether prediction or fallback occurred. Record commands, errors, task completion and feedback. Present native-versus-Tandem output and timing evidence at the demonstration and ask whether the observed tradeoff meets the intended serving use. This activity is planned; no stakeholder participation or approval is assumed.

\section{System Tests}
\subsection{Common Fixtures and Execution Rules}
Before execution, freeze exact source/target model revisions, tokenizer/chat formats, predictor checkpoint and alignment, precision, engine/dependency versions, hardware, resource limits and seeds in a run manifest. Use public or synthetic conversations split by conversation group into fitting, development and final evaluation sets. Related prefixes/branches from one conversation must stay in the same split. Record dataset provenance and split IDs; do not fit or tune on final evaluation.

Define fixture F1 as a two-turn synthetic conversation: system text ``Return only the code supplied in the notes'', shared history containing ``The code is EQBU'', and latest user text ``What is the code?''. Retain full rendered prompts and target token positions in explicitly enabled synthetic test traces. Native output is the behavioural oracle; EQBU is a task-success observation, not a guaranteed model response. Expand with fixed filler to short, medium and supported-limit histories, placing the code early, middle and late. Exact token counts, dataset size and supported model pair remain open decisions and must be fixed before execution.

For each automatic case, start from a recorded clean configuration and capture output events, final status, reuse reason and relevant counters. Run positive controls as well as individual negative mutations. Correctness traces use synthetic/public data and a declared retention period; normal exports follow the privacy cases below.

\subsection{Tests for Functional Requirements}
\case{F01}{Request validation}{FR-01, FR-02}
\textbf{Control:} Automatic. \textbf{Initial state:} Healthy engine, prediction disabled.
\textbf{Input:} F1 with a known target; then separately remove a required field, use an unknown model, duplicate a request ID, set output cap to zero and exceed the declared token limit by one.
\textbf{Expected:} The valid request preserves model, ordered prompt and settings. Invalid requests are rejected with reasons before generation or associated oversized allocation.
\textbf{Procedure/oracle:} Inspect admission events, native-rendered tokens and allocation counters against the SRS Data Dictionary.

\case{F02}{Prefix and cache eligibility}{FR-03, FR-04, SR-02}
\textbf{Control:} Automatic. \textbf{Initial state:} Qualified setup and valid F1 prediction.
\textbf{Input:} Positive control followed by independent changes to message order, system text, prefix positions, tokenizer/format, model/predictor/configuration versions, tensor shape, precision, readiness and finite values (NaN/infinity).
\textbf{Expected:} Valid predictions are eligible; every incompatible or incomplete cache is blocked from loading. A valid request with an unusable optional cache follows native processing, or reports a resource/server failure if native processing is unavailable.
\textbf{Procedure/oracle:} Compare the validation decision and actual load events; equal tensor dimensions alone must not permit reuse.

\case{F03}{Predict, load and prefill the suffix}{FR-05, FR-16}
\textbf{Control:} Automatic GPU test. \textbf{Initial state:} Frozen base models, versioned predictor and held-out F1 variants.
\textbf{Input:} Source KV with declared alignment and an uncached latest turn.
\textbf{Expected:} Prediction uses source KV without constructing it through target-prefix prefill; accepted state covers the intended prefix and the suffix is processed once in order. Output passes A01.
\textbf{Procedure/oracle:} Instrument actual target token work and load positions. Compute native target references offline, separately from the timed prediction path, and apply A01. Save architecture, fitting and parameter versions.

\case{F04}{Native path and recovery}{FR-06, FR-07, NFR-06}
\textbf{Control:} Automatic fault injection. \textbf{Initial state:} Healthy engine.
\textbf{Input:} F1 with prediction disabled, no source cache, unsupported pair, same-model turn and unshared history. Separately inject predictor and loader exceptions before output, then an execution failure after the first output event.
\textbf{Expected:} Ordinary cases use native processing with native prefix caching retained. Recoverable pre-output failures permit at most one native restart under the same ID; a failed restart or post-output failure ends as failed without a mixed or duplicate answer. A later valid request completes if the engine is healthy.
\textbf{Procedure/oracle:} Compare native output, restart counts, event order and final statuses to the SRS recovery rules.

\case{F05}{Concurrent lifecycle and cleanup}{FR-08, FR-09, FR-15, SR-03}
\textbf{Control:} Automatic. \textbf{Initial state:} Two active requests sharing permitted cache state.
\textbf{Input:} Complete one request while the other decodes; repeat with queued/running cancellation, timeout, forced failure and attempted cache eviction/replacement. Repeat for $N_{cycle}$ cycles.
\textbf{Expected:} Events remain ordered under the correct ID; each accepted request has exactly one final status. Ended requests retain no private allocation, and active requests retain needed cache state and pass A01. Policy-retained reusable caches are accounted for separately.
\textbf{Procedure/oracle:} Inspect allocation ownership and reference counts as well as measured memory; allocator-reserved memory alone is not a leak oracle. Check safe stopping points and timeout reasons using $T_{timeout}$.

\case{F06}{Configuration and qualification}{FR-13, FR-14, SR-05}
\textbf{Control:} Automatic plus evidence review. \textbf{Initial state:} Unqualified setup.
\textbf{Input:} Attempt enablement with absent/failed accuracy, isolation or recovery evidence; then with complete passing evidence. Change a model/predictor version and restart the server while retaining an old handle.
\textbf{Expected:} Missing or failed evidence prevents approval; approved records identify splits, parameters and limits. Old/incompatible handles are blocked and untested combinations use native processing. Invalid settings are rejected.
\textbf{Procedure/oracle:} Inspect qualification records and verify actual reuse decisions before and after configuration changes.

\case{F07}{Measurement and replay integrity}{FR-10, FR-11, FR-12, NFR-01, NFR-12, DR-02}
\textbf{Control:} Automatic schema checks and manual report review.
\textbf{Input:} A fixed batch containing success, fallback, rejection, cancellation and failure, replayed natively and with Tandem.
\textbf{Expected:} IDs, versions, settings, arrivals and limits match; unexplained differences invalidate comparison. Reports preserve reasons, counts, units, actual prompt work, stage timings, output differences, memory and GPU-time measurements. Missing measurements have reasons rather than fabricated estimates. Codes match documented corrective actions.
\textbf{Procedure/oracle:} Reconcile request/event records with the report, deliberately alter one replay setting to test mismatch detection, and verify that timed totals include prediction, validation and loading. Record offline development costs separately and document assumptions for any priced estimate.

\subsection{Tests for Nonfunctional Requirements}
\case{A01}{KV and output accuracy}{NFR-05, FR-14, FR-16}
\textbf{Control:} Automatic GPU comparison. \textbf{Input:} Held-out fixtures for each supported setup, including declared context limits.
\textbf{Procedure:} First establish native repeatability. For keys and values separately, report $\|\widehat{K}-K\|_F$ and its ratio to $\|K\|_F$ per layer/cache group (likewise for values). A zero reference norm gives undefined relative error and a reported absolute error. Do not invent a KV-error pass threshold.

Compare corresponding logits after identical preceding tokens, with identical masks. Require $|z_T-z_N|\leq\epsilon_a+\epsilon_r|z_N|$ for every compared finite score; unexpected nonfinite values fail. Then independently run greedy decoding and require identical tokens for $N_{decode}$ tokens or an earlier shared stop. Report first divergence and aggregate failures. Native instability blocks a reliable qualification conclusion. Sampling needs separate recorded qualification or native fallback. These draft limits require OI3 review before final evaluation.

\case{A02}{Privacy and text preservation}{SR-01, SR-04, SR-06, NFR-03}
\textbf{Control:} Automatic. \textbf{Input:} F1 across different host-granted sharing groups; forged request-side permissions; concurrent branches whose private suffix alternates between ``secret-A'' and ``secret-B''. Add accents, non-Latin text, whitespace and delimiter variants.
\textbf{Expected:} Unauthorized cache state is never used or exposed. Holding the victim request fixed, changing another branch's private content preserves victim logits/tokens under A01 with passing native controls. Rendered prompt tokens match native processing. Normal logs and exports contain no planted secrets, raw prompts, answers, KV or credentials.
\textbf{Procedure/oracle:} Inspect cache access decisions, compare victim runs and scan normal exports using unique synthetic markers, including markers supplied in request-controlled fields.

\case{P01}{Latency, throughput and ordinary-serving overhead}{NFR-04, NFR-07, NFR-08}
\textbf{Control:} Automatic GPU benchmark. \textbf{Input:} Preselected model-switch/shared-history workloads and separate same-model, unshared, disabled-prediction and missing-cache controls. Cover cold/warm caches, supported lengths and concurrency. Freeze workloads before final runs.

\textbf{Procedure:} Keep native prefix caching equally enabled. Match arrivals, routing, generation settings and resource limits. Record source-cache preparation and model-loading policy. Run $N_{warmup}$ untimed warmups and $N_{pairs}$ alternating timed pairs; use the same declared cache preparation for each arm. Timing includes queueing, retrieval, prediction, validation/loading, suffix prefill and decoding. Report median/p95 TTFT, end-to-end time, successful requests per second, prompt work, total GPU time and GPU time per completion. TTFT is missing for requests returning no tokens. Throughput spans first arrival to last final status and counts successful completions only.

Bootstrap complete paired trials $N_{boot}$ times with a recorded seed to obtain $C$ confidence intervals; individual requests in one trial are not independent repetitions. For positive native metrics, calculate percentage change relative to native. NFR-04 needs fewer computed prompt tokens and a TTFT-change interval wholly below zero on at least one preselected workload, with passing output checks and no extra failures. NFR-07 needs a throughput-change interval wholly above zero and an upper p95 TTFT-change bound no greater than $\delta$. For NFR-08, upper latency-change bounds must be at most $\delta$ and the lower throughput-change bound at least $-\delta$. Report every workload, including regressions and inconclusive intervals. Any additional-trial policy must be recorded before final evaluation; avoid repeatedly sampling merely until a favourable interval appears.

\case{D01}{Installation, supported limits and reproduction}{NFR-02, NFR-09, NFR-10, NFR-12, DR-01}
\textbf{Control:} Manual task-based inspection with automated boundary checks.
\textbf{Input:} Clean supported environment, release instructions and manifest; token/concurrency values at and just beyond published limits.
\textbf{Expected:} An independent reviewer can install, run a native and prediction example, interpret fallback and reproduce basic/output checks. Boundary behaviour agrees with the documented support limits; unsupported setups are not presented as qualified.
\textbf{Procedure:} Record all commands, deviations and errors. Compare hardware/load differences before interpreting repeated timing results. File unclear instructions and missing status-code explanations as issues.

\subsection{Traceability Between Test Cases and Requirements}
This initial mapping covers the numbered FR, NFR, SR and documentation requirements used above. Review it against the approved SRS after scope changes; a mapped case is not evidence of a passing result.
\begin{longtable}{@{}p{0.58\textwidth}p{0.34\textwidth}@{}}
\toprule Requirements & Cases\\\midrule\endhead
FR-01, FR-02 & F01\\
FR-03, FR-04, SR-02 & F02\\
FR-05, FR-16 & F03, A01\\
FR-06, FR-07, NFR-06 & F04\\
FR-08, FR-09, FR-15, SR-03 & F05\\
FR-13, SR-05 & F06\\
FR-14 & F06, A01, A02, F04\\
FR-10, FR-11, FR-12, NFR-01, DR-02 & F07\\
NFR-05 & A01\\
NFR-03, SR-01, SR-04, SR-06 & A02\\
NFR-04, NFR-07, NFR-08 & P01\\
NFR-02, NFR-09, NFR-10, DR-01 & D01\\
NFR-12 & F07, D01\\\bottomrule
\end{longtable}

\section{Unit Test Description}
Detailed module cases and module traceability are deferred until the MIS defines interfaces and state ownership. The following strategy guides that work without inventing completed design decisions.
\subsection{Unit Testing Scope}
Prioritize Tandem-owned prefix/alignment checks, predictor adapters, cache validation/loading, lifecycle handling and measurement calculations. Use small synthetic tensors and fake engine handles for CPU tests; GPU integration remains necessary for actual loading and generation. External model internals are outside unit scope.
\subsection{Tests for Functional Requirements}
For each MIS access program, cover normal inputs, boundaries and each documented exception. Candidate checks include zero/nonfinite tensor handling, off-by-one suffix positions, stale versions, denied sharing, one-restart enforcement and final-status uniqueness. Deterministic fixtures must specify expected outputs independently of the implementation.
\subsection{Tests for Nonfunctional Requirements}
Test error-norm calculations against hand-computed tensors, including zero references; measurement aggregation against a small known event stream; and missing TTFT handling when no token is returned. GPU predictor/loading microbenchmarks explain stage overhead but do not replace P01.
\subsection{Traceability Between Test Cases and Modules}
Pending the MIS: assign each module its relevant system requirements, unit-test IDs and an owner, and document any excluded module with a rationale.

\section{Appendix}
\subsection{Symbolic Parameters}
These values are inherited from the SRS draft, not newly approved thresholds. OI3 must be resolved before final tests.

\noindent\begin{tabularx}{\textwidth}{llX}
\toprule Symbol & Draft value & Purpose\\\midrule
$\epsilon_a$ & 0.001 & Absolute logit tolerance.\\
$\epsilon_r$ & 0.001 & Relative logit tolerance.\\
$N_{decode}$ & 32 & Greedy comparison length.\\
$N_{cycle}$ & 100 & Cleanup cycles.\\
$T_{timeout}$ & 120 seconds & Default request timeout.\\
$N_{warmup}$ & 3 & Untimed warmups.\\
$N_{pairs}$ & 10 & Initial timed pairs.\\
$N_{boot}$ & 2,000 & Bootstrap resamples.\\
$C$ & 95\% & Confidence level.\\
$\delta$ & 5\% & Allowed slowdown/loss.\\\bottomrule
\end{tabularx}

\subsection{Open Decisions and Next Steps}
Before running this plan, resolve SRS OI1--OI3 and OI7: choose the model pair, predictor/alignment, dataset/split sizes, exact versions, supported token/concurrency limits and GPU budget; review draft tolerances on development data. Record seeds, additional-trial policy, trace retention and evidence locations. Assign named test owners and an independent reviewer. Align earlier scope documents under OI5/OI6 and confirm custom-document approval. First implement native controls and A01, then prediction/loading and failure tests; run performance qualification only after correctness and isolation pass.

\subsection{Usability Inspection Questions}
During D01 ask: Which setup is supported? How do you enable or disable prediction? Did this request reuse KV or fall back, and why? Where are the limitations and reproduction commands? Record misunderstandings and revise the instructions. No general end-user usability survey is planned for this server-side prototype.

\section{Appendix --- Generative AI Use Disclosure}
\subsection{AI Tools Used}
OpenAI Codex assisted with this initial draft.
\subsection{How and Where Used}
Codex replaced template instructions with proposed scope, review processes, system cases, traceability, unit-test strategy and appendices using repository documents and the user's clarification that Tandem performs cross-model KV prediction.
\subsection{Verification of AI-Generated Content}
The draft was cross-checked against the Volere SRS and repository checklist. Team review of feasibility, requirement coverage and proposed procedures remains pending. No experiments, stakeholder validation or qualification results were generated by this drafting work. The team remains responsible for the submitted document.

\section{Appendix --- Reflection}
Pending team contributions after drafting and review: record what went well, difficulties and resolutions, and each member's required learning. For each learning area, identify two acquisition approaches and the member's chosen approach. Do not substitute proposed plans for personal experiences.
\end{document}
